\chapter{L'oracolo seriale}
\label{app:oracolo}

La traccia richiede il collaudo funzionale del prodotto distribuito contro
un'implementazione seriale minimale. Il progetto la realizza in
\file{src/serial/serial.c} e la usa come \emph{oracolo}: non è un backend da
misurare, è il termine di paragone rispetto al quale ogni altro percorso di
calcolo viene dichiarato corretto.

\section{L'implementazione}
\label{sec:oracolo-implementazione}

\code{serial\_gemm} è la trascrizione letterale della \cref{eq:definizione}:
tre cicli annidati, nessuna ottimizzazione, nessuna dipendenza da MPI o da CUDA.

\begin{lstlisting}[language=C,caption={Il nucleo di \file{src/serial/serial.c}.},label={lst:oracolo}]
for (i = 0; i < m; i++) {
    for (c = 0; c < k; c++) {
        scalar_t s = (scalar_t)0;
        for (j = 0; j < n; j++)
            s += A[(size_t)i * lda + j] * X[(size_t)j * ldx + c];
        Y[(size_t)i * ldy + c] = s;
    }
}
\end{lstlisting}

Due scelte meritano una nota.

L'ordine dei cicli è $i \to c \to j$, cioè lo schema~B della
\cref{sec:schemi}, mentre tutti i backend misurati adottano lo schema~A. La
scelta è deliberata e non ha nulla a che vedere con le prestazioni: se il
ragionamento sull'ordinamento dei cicli fosse sbagliato, un oracolo costruito
con la stessa struttura dei backend replicherebbe l'errore invece di
rivelarlo. L'oracolo deve essere \emph{ovviamente} corretto, non veloce.

La firma accetta le tre leading dimension $\code{lda}$, $\code{ldx}$ e
$\code{ldy}$ separatamente dalle dimensioni logiche: l'oracolo può quindi
operare su buffer con padding esattamente come i backend, e la validazione
resta valida anche nelle build che forzano $\code{lda} \neq \nloc$.
Il tipo \code{scalar\_t} è lo stesso della build corrente
(\cref{sec:precisione}): l'oracolo viene compilato nella precisione che si sta
validando, non in doppia precisione per confronto.

\section{Il confronto}
\label{sec:oracolo-confronto}

Con l'opzione \code{--check} il driver esegue, dopo le
ripetizioni cronometrate e fuori da qualunque cronometro, la procedura
\code{check\_against\_serial}:

\begin{enumerate}[nosep,left=0pt]
  \item i processi della prima colonna della griglia - l'unica che al termine
        della riduzione possiede $Y$ per intero - raccolgono i propri blocchi
        di righe sul rank~0 con una \code{MPI\_Gatherv} sul comunicatore di
        colonna;
  \item il rank~0 \emph{ricostruisce da sé} $A$ e $X$ globali invocando il
        generatore deterministico della \cref{eq:gen} con lo stesso seed, senza
        riceverli da nessun processo;
  \item calcola $Y_{\text{ref}}$ con \code{serial\_gemm} e confronta.
\end{enumerate}

Il passo~2 è ciò che rende il collaudo indipendente dal codice che deve
collaudare: se l'oracolo ricevesse $A$ e $X$ dagli altri processi, un errore
nella distribuzione dei dati si propagherebbe identico nei due termini del
confronto e il test passerebbe. Ricostruendoli dal generatore, l'oracolo ha
accesso al dato globale \emph{per definizione} e non al dato distribuito.

La metrica del confronto è l'errore relativo in norma $L_2$ sull'intera $Y$,
\begin{equation}
  \mathrm{err}_{\text{rel}}
  =
  \frac{\lVert Y - Y_{\text{ref}} \rVert_2}{\lVert Y_{\text{ref}} \rVert_2},
  \label{eq:err-rel-oracolo}
\end{equation}
riportata nella colonna \code{rel\_err} del CSV e posta a $-1$ quando
\code{--check} non è attivo.

\section{La tolleranza}
\label{sec:oracolo-tolleranza}

La soglia del confronto non può essere una costante. Ogni elemento di $Y$ è una
riduzione su $N$ termini, quindi l'errore di arrotondamento accumulato dipende
da $N$; e poiché il generatore produce valori a media nulla, gli arrotondamenti
hanno segno indipendente e l'errore cresce come $\varepsilon\sqrt{N}$, non come
il pessimistico $\varepsilon N$ del caso in cui tutti i termini abbiano lo
stesso segno. La soglia adottata segue quindi la stessa legge, con un fattore
di sicurezza:
\begin{equation}
  \mathrm{tol}(N) = S \cdot \varepsilon \cdot \sqrt{N},
  \qquad S = 8,
  \label{eq:tol-oracolo}
\end{equation}
dove $\varepsilon$ è l'epsilon di macchina del tipo scalare della build
(\code{DBL\_EPSILON} oppure \code{FLT\_EPSILON}). La formula è implementata in
\code{scalar\_check\_tol} in \file{src/common/scalar.h}, cioè accanto alla
definizione di \code{scalar\_t}: precisione e tolleranza cambiano insieme e non
possono divergere.

Il valore di $N$ che entra nella \cref{eq:tol-oracolo} è il numero
\emph{globale} di colonne di $A$, non $\nloc$: è su tutte e $N$ le colonne che
l'oracolo somma, indipendentemente da come il lavoro sia stato distribuito.

L'esito è binario. Il driver stampa \code{PASS} o \code{FAIL} e, soprattutto,
termina con stato di uscita non nullo in caso di \code{FAIL}: un errore di
correttezza interrompe lo script della campagna invece di finire silenziosamente
in una riga di CSV.

\section{Perché resta fuori dalle misure}
\label{sec:oracolo-costo}

L'oracolo costa $O(MNk)$ operazioni seriali su un solo processo e richiede al
rank~0 la memorizzazione di $A$ e $X$ globali, cioè $O(MN)$ elementi: alle
taglie delle campagne è più costoso, da solo, dell'intera sessione di misura.
Per questo la validazione è un'esecuzione a sé, condotta su taglie ridotte
\emph{prima} della campagna corrispondente, e \code{--check} non compare mai
nei comandi che producono i tempi riportati nel \cref{cap:risultati}.

La validazione va inoltre ripetuta per ogni combinazione di backend e
precisione: l'ordine della riduzione cambia da un kernel all'altro, e un
backend validato solo in doppia precisione non è validato.
